% ============================================================
%  part4/ch21-energy-dissipation.tex
%  Chapter 21: Energy Dissipation Theorems
% ============================================================

\chapter{Energy Dissipation Theorems}
\label{ch:energy-dissipation}

\section{Advective RSVP and transport conservation}

The more physical RSVP system includes transport by $\mathbf{v}$:
\begin{align*}
  \partial_t\Phi + \mathbf{v}\cdot\nabla\Phi
  &= -M_\Phi\frac{\delta\mathcal{F}}{\delta\Phi},\\
  \partial_t S + \mathbf{v}\cdot\nabla S
  &= -M_S\frac{\delta\mathcal{F}}{\delta S} + \sigma,\\
  \partial_t L + \nabla\cdot(L\mathbf{v})
  &= \nabla\cdot(M_L\nabla\mu_L).
\end{align*}

\begin{proposition}[Advection is energy-neutral]
  \label{prop:adv-neutral}
  \statusproven{}
  For incompressible flow $\nabla\cdot\mathbf{v} = 0$ and
  periodic or no-flux boundaries,
  pure advection $\partial_t\Phi = -\mathbf{v}\cdot\nabla\Phi$
  satisfies $\frac{d}{dt}\int\frac{1}{2}\Phi^2\,dx = 0$.
\end{proposition}

\begin{proof}
  $\frac{d}{dt}\int\frac{1}{2}\Phi^2\,dx
  = \int\Phi\partial_t\Phi\,dx
  = -\int\Phi\mathbf{v}\cdot\nabla\Phi\,dx
  = -\int\mathbf{v}\cdot\nabla(\frac{1}{2}\Phi^2)\,dx$.
  By the divergence theorem and $\nabla\cdot\mathbf{v}=0$:
  $= \int\frac{1}{2}\Phi^2\nabla\cdot\mathbf{v}\,dx = 0$.
\end{proof}

\begin{theorem}[Energy dissipation with incompressible transport]
  \label{thm:rsvp-adv-dissipation}
  \statusproven{}
  Assume periodic or no-flux boundaries and $\nabla\cdot\mathbf{v}=0$.
  If advective terms are skew-adjoint with respect to the
  energy pairing and the entropy source satisfies
  $\int\frac{\delta\mathcal{F}}{\delta S}\sigma\,dx \leq 0$,
  then $\frac{d\mathcal{F}}{dt} \leq 0$.
\end{theorem}

\begin{proof}
  Decompose $\frac{d\mathcal{F}}{dt} = I_{\mathrm{adv}} + I_{\mathrm{diss}}$.
  By Proposition~\ref{prop:adv-neutral} and skew-adjointness,
  $I_{\mathrm{adv}} = 0$.
  The dissipative part satisfies
  $I_{\mathrm{diss}} \leq 0$ by the gradient-flow structure
  and the entropy-source condition.
\end{proof}

\section{Entropy smoothing}

\begin{proposition}[Entropy gradient decay]
  \label{prop:entropy-smooth}
  \statusproven{}
  If $\partial_t S = \kappa\Delta S$ with $\kappa > 0$, then
  \[
    \frac{d}{dt}\int_\Omega \tfrac{1}{2}|\nabla S|^2\,dx
    = -\kappa\int_\Omega(\Delta S)^2\,dx \leq 0.
  \]
\end{proposition}

\begin{proof}
  $\frac{d}{dt}\int\frac{1}{2}|\nabla S|^2\,dx
  = \int\nabla S\cdot\nabla(\partial_t S)\,dx
  = \kappa\int\nabla S\cdot\nabla\Delta S\,dx
  = -\kappa\int(\Delta S)^2\,dx \leq 0$.
\end{proof}

\section{Lamphrodyne as entropy relaxation pressure}

\begin{definition}[Lamphrodyne]
  \label{def:lamphrodyne}
  \[
    D = -\frac{\delta\mathcal{F}}{\delta S}
    = \chi\Delta S - Q'(S) + T.
  \]
\end{definition}

Entropy relaxation takes the Cahn--Hilliard-like form:
$\partial_t S = \nabla\cdot(M_S\nabla D) + \sigma$.

\begin{StatusBox}{Proven Framework}
  Lamphrodyne $D$ is not an independent field.
  It is the negative variational derivative of $\mathcal{F}$
  with respect to $S$.
  Its derivation from $\mathcal{F}$ is proven.
  Its cosmological interpretation as void-expansion pressure
  is motivated (Chapter~\ref{ch:void-formation}) but requires
  eikonal derivation (Chapter~\ref{ch:redshift}).
\end{StatusBox}
